\chapter*{Lampiran}
\addcontentsline{toc}{section}{A. Struktur Repositori}

Struktur repositori yang dipakai selama pengerjaan ditunjukkan pada Tabel \ref{tab:struktur}. Berkas data, checkpoint, dan hasil eksperimen tidak dimasukkan ke dalam repositori karena ditangani melalui \texttt{.gitignore}.

\begin{table}[H]
\centering
\caption{Struktur direktori repositori}
\label{tab:struktur}
\begin{tabular}{|L{3.9cm}|L{7.1cm}|}
\hline
\textbf{Path} & \textbf{Keterangan} \\
\hline
\texttt{src/fashion\_cgan/models.py} & Generator dan Discriminator MLP beserta validasi masukan \\
\hline
\texttt{src/fashion\_cgan/data.py} & Loader Fashion-MNIST dengan normalisasi $[-1,1]$ dan seed \\
\hline
\texttt{src/fashion\_cgan/train.py} & Prosedur pelatihan, checkpoint, metrik, opsi resume \\
\hline
\texttt{src/fashion\_cgan/generate.py} & Pembangkitan grid citra dari checkpoint \\
\hline
\texttt{src/fashion\_cgan/evaluate.py} & Classifier evaluasi dan pengukuran akurasi per kelas \\
\hline
\texttt{src/fashion\_cgan/utils.py} & Fungsi bantu denormalisasi, metrik, dan penyimpanan grid \\
\hline
\texttt{configs/default.yaml} & Konfigurasi eksperimen \\
\hline
\texttt{tests/} & Pengujian unit model, data, training, generate, evaluasi \\
\hline
\texttt{notebooks/} & Notebook visualisasi laporan dan bahan presentasi \\
\hline
\texttt{docs/} & Arsitektur, rencana implementasi, evaluasi, log sesi AI \\
\hline
\texttt{evidence/} & Bukti sesi pengerjaan tiap anggota \\
\hline
\end{tabular}
\end{table}

\addcontentsline{toc}{section}{B. Perintah Reproduksi}

Seluruh tahapan dapat dijalankan ulang dengan perintah pada Tabel \ref{tab:perintah}. Diasumsikan lingkungan virtual pada direktori \texttt{.venv} sudah terpasang melalui \texttt{pip install -r requirements.txt}.

\begin{table}[H]
\centering
\caption{Perintah reproduksi eksperimen}
\label{tab:perintah}
\begin{tabular}{|L{3.4cm}|L{7.9cm}|}
\hline
\textbf{Tahap} & \textbf{Perintah} \\
\hline
Pengujian & \texttt{python -m pytest tests -q} \\
\hline
Training & \texttt{python -m fashion\_cgan.train --device cpu} \\
\hline
Resume & \texttt{python -m fashion\_cgan.train --resume outputs/cgan\_epoch\_N.pt} \\
\hline
Pembangkitan & \texttt{python -m fashion\_cgan.generate outputs/cgan\_epoch\_50.pt} \\
\hline
Evaluasi & \texttt{python -m fashion\_cgan.evaluate outputs/cgan\_epoch\_50.pt} \\
\hline
Visualisasi & Eksekusi \texttt{notebooks/report\_visuals.ipynb} dari root repositori \\
\hline
\end{tabular}
\end{table}

Pada Windows PowerShell, variabel \texttt{PYTHONPATH} perlu diisi dengan nilai \texttt{src} agar paket \texttt{fashion\_cgan} dapat diimpor.

\addcontentsline{toc}{section}{C. Bukti Sesi Pengerjaan}

Bukti pemakaian asisten AI disimpan terpisah per anggota pada direktori \texttt{evidence/}. Ringkasannya disajikan pada Tabel \ref{tab:evidence}.

\begin{table}[H]
\centering
\caption{Ringkasan berkas bukti per anggota}
\label{tab:evidence}
\begin{tabular}{|L{3.9cm}|L{7.1cm}|}
\hline
\textbf{Anggota} & \textbf{Berkas bukti} \\
\hline
Ivander Dzaky Khairullah & \texttt{session-001-planning.md}, \texttt{session-002-generator-discriminator.md}, \texttt{session-004-preprocessing-training.md} beserta ekspor JSON asli \\
\hline
Muhammad Ragil & \texttt{session-001-data-training-and-github-workflow.md} \\
\hline
Silmi Sila Nia Uma & \texttt{session-001-evaluation.md} \\
\hline
\end{tabular}
\end{table}
